\documentclass[12pt,a4paper]{article}

\usepackage[margin=2.5cm]{geometry}
\usepackage{setspace}
\usepackage{parskip} % space between paragraphs, no indentation

\begin{document}
	
	\begin{center}
		\Large Self‑Reflection on Thesis Writing
	\end{center}
	
	\onehalfspacing
	
	% Start writing your 500-word essay here:
	In this reflection, I evaluate my thesis writing experience, focusing on my attitudes, behaviors, and motivation during the writing process. My thesis investigates how effectively machine learning models can assess pilots’ performance in virtual reality flight tasks using multiple inputs. One of the most challenging aspects was reviewing and synthesizing previous studies to support and justify my reasoning, which demanded substantial time and mental effort.

	However, it was also exciting and interesting to explore current technological developments in the field. At the same time, it was a valuable experience to learn how researchers design experiments and conduct analyses. This kept me motivated to formulate and refine my choices of relevant elements in the machine learning models and in the corresponding written justifications.
	
	The writing of the methodology and results sections of the thesis also proved to be challenging. It was more manageable to implement the methods in data preprocessing and model setup than to translate the actual implementation into coherent and clear written text. To provide a precise and accurate description, I carefully considered multiple wording options and made sure that the accompanying tables and figures helped to illustrate and clarify the intended meaning. The discussion section also required a great deal of time and effort to finalize. Following the same principles as in the preceding sections, I checked the statistical results and reasoning multiple times to ensure that the interpretation was logical and accurately reflected the outputs.
	
	Technical and time management issues did not arise during the writing process. However, there are several improvements that can be applied to future writing activities. First, more time should be spent practicing academic writing, as this will help generate concise and clear sentences, avoid vagueness, and save time by conveying meaning without redundant expressions. Second, the writing process can run in parallel with the experimental procedure, making it much easier to correct and adjust the text while the understanding of the outcomes is still fresh. Finally, more communication should take place between group members to proofread each other’s work and to ensure that the thesis is generally understandable to readers with a similar background.
	
	Overall, the thesis writing process has increased my awareness of both my strengths and weaknesses in academic writing. I discovered that I am persistent and careful when checking results and justifying methodological choices, but that I still need more practice in expressing complex ideas clearly and efficiently. This experience has also given me a more realistic view of research, showing me that confusion and revision are natural parts of the process. In future projects, I intend to apply these lessons by starting the writing earlier, working more closely with feedback, and continuously reflecting on how my attitudes and behaviors influence the quality of my work.
	
	
\end{document}